\documentclass[11pt,a4paper]{article}
\usepackage[italian]{babel}
\usepackage[margin=2.5cm]{geometry}
\usepackage{parskip}
\usepackage{amsmath, amssymb, amsthm}
\usepackage{booktabs}
\usepackage{array}
\usepackage{xcolor}
\usepackage{needspace}
\usepackage{enumitem}
\usepackage{listings}
\usepackage{tikz}
\usepackage[most]{tcolorbox}
\usepackage[hidelinks]{hyperref}
\usetikzlibrary{decorations.pathreplacing, shapes.geometric, arrows.meta, positioning}

% Niente vedove/orfane: i paragrafi non lasciano righe isolate a fine/inizio pagina
\widowpenalty=10000
\clubpenalty=10000

% Elenchi più compatti
\setlist{itemsep=2pt, parsep=0pt, topsep=4pt, partopsep=0pt}

% --- Riquadri con bordo nero, senza numero (si spezzano tra due pagine) ---
% Uso: \begin{definizione} ... \end{definizione}
%      \begin{teorema}[Nome] ... \end{teorema}  (il nome è facoltativo)
\tcbset{nota/.style={enhanced jigsaw, breakable, colback=white,
  colframe=black, boxrule=0.5pt, sharp corners}}
\NewTColorBox{definizione}{}{nota,
  before upper={\textbf{Definizione.}\ }}
\NewTColorBox{teorema}{o}{nota,
  before upper={\textbf{Teorema\IfValueT{#1}{ (#1)}.}\ }}
\NewTColorBox{proposizione}{o}{nota,
  before upper={\textbf{Proposizione\IfValueT{#1}{ (#1)}.}\ }}
\NewTColorBox{proprieta}{o}{nota,
  before upper={\textbf{Proprietà\IfValueT{#1}{ (#1)}.}\ }}
\NewTColorBox{assioma}{o}{nota, boxrule=1pt,
  before upper={\textbf{Assioma\IfValueT{#1}{ (#1)}.}\ }}

% --- Ambienti semplici (senza riquadro) ---
% Il titolo sta in cima al blocco, anche se il contenuto inizia con un riquadro o
% con codice; e non resta mai da solo in fondo a una pagina.
% Uso: \begin{esempio}[Titolo facoltativo] ... \end{esempio}
\newcommand{\newrunin}[2]{%
  \NewDocumentEnvironment{#1}{o}{%
    \par\Needspace{4\baselineskip}\addvspace{\medskipamount}\noindent
    \textbf{#2}\IfValueT{##1}{ (##1)}\textbf{.}\ \ignorespaces}%
  {\par\addvspace{\medskipamount}}}
\newrunin{esempio}{Esempio}
\newrunin{esercizio}{Esercizio}
\newrunin{osservazione}{Osservazione}

% V e F nelle tavole di verità
\newcommand{\V}{\textsf{V}}
\newcommand{\F}{\textsf{F}}

% --- Codice C ---
% Uso: \begin{codice} ... \end{codice}      codice C numerato
%      \begin{comando} ... \end{comando}    comandi da terminale
%      \begin{risultato} ... \end{risultato} output di un programma
%      \lstinline|printf("ciao");|           codice dentro al testo
\lstdefinestyle{c}{
  language=C,
  basicstyle=\ttfamily\small,
  keywordstyle=\color{blue}\bfseries,
  commentstyle=\color{gray}\itshape,
  stringstyle=\color{red!70!black},
  numbers=left,
  numberstyle=\tiny\color{gray},
  numbersep=8pt,
  columns=flexible,
  keepspaces=true,
  breaklines=true,
  showstringspaces=false,
  tabsize=4,
  morekeywords={include, define},
  % lettere accentate nei commenti
  literate={à}{{\`a}}1 {è}{{\`e}}1 {é}{{\'e}}1 {ì}{{\`i}}1
           {ò}{{\`o}}1 {ù}{{\`u}}1 {È}{{\`E}}1 {À}{{\`A}}1
}
\lstset{style=c}
\newtcblisting{codice}{listing only, nota, left=6mm,
  listing options={style=c}}
\newtcblisting{comando}{listing only, nota,
  listing options={basicstyle=\ttfamily\small, numbers=none, language={},
    columns=flexible, keepspaces=true, breaklines=true}}
\newtcblisting{risultato}{listing only, enhanced jigsaw, breakable,
  colback=black!5, colframe=black, boxrule=0.5pt, sharp corners,
  listing options={basicstyle=\ttfamily\small, numbers=none, language={},
    columns=flexible, keepspaces=true, breaklines=true}}

% --- Diagrammi di flusso (simboli standard) ---
\tikzset{
  terminale/.style = {ellipse, draw, minimum width=2.5cm, minimum height=0.9cm},
  io/.style        = {trapezium, trapezium left angle=70, trapezium right angle=110,
                      draw, minimum width=2.5cm, minimum height=0.9cm},
  processo/.style  = {rectangle, draw, minimum width=2.5cm, minimum height=0.9cm},
  decisione/.style = {diamond, draw, aspect=2, minimum width=2.5cm},
  freccia/.style   = {thick, -{Stealth}}
}

\title{MCD, mcm e algoritmo di Euclide}
\author{Marco Cerbella}
\date{Lezione 11 -- 26 ottobre 2026}

\begin{document}
\maketitle

\section{Massimo comune divisore}

\begin{definizione}
Siano $a, b \in \mathbb{Z}$ non entrambi nulli. Un elemento $d \in \mathbb{Z}$ è \emph{un massimo comun divisore} tra $a$ e $b$ se
\begin{enumerate}
  \item[i)] $d \mid a$ e $d \mid b$;
  \item[ii)] $\forall d' \in \mathbb{Z}$ tale che $d' \mid a$ e $d' \mid b$ si ha $d' \mid d$.
\end{enumerate}
\end{definizione}

\begin{osservazione}
In generale, se $d$ è un massimo comun divisore tra $a$ e $b$, anche $-d$ è un massimo comun divisore tra $a$ e $b$. Chiamiamo $\operatorname{MCD}(a, b)$ quello \emph{positivo}.
\end{osservazione}

\begin{esempio}
$\operatorname{MCD}(48, 15) = 3$, \quad $\operatorname{MCD}(3, 7) = 1$, \quad $\operatorname{MCD}(-12, 36) = 12$.
\end{esempio}

\begin{proprieta}
\begin{itemize}
  \item $\operatorname{MCD}(a, b) = \operatorname{MCD}(b, a) = \operatorname{MCD}(|a|, |b|)$.
  \item $\operatorname{MCD}(ab, ac) = |a|\operatorname{MCD}(b, c)$.
  \item $\operatorname{MCD}(0, a) = |a|$: infatti $d \mid 0$ per ogni $d$ (ogni $d$ divide $0$), e $d \mid a$; quindi $d$ è un divisore di $a$ e il più grande possibile è $|a|$.
  \item $\operatorname{MCD}(0, 0)$ \emph{non è definito}.
\end{itemize}
\end{proprieta}

\begin{definizione}[Numeri coprimi]
Due numeri $a, b$ sono \emph{coprimi} (o \emph{relativamente primi}, o \emph{primi tra loro}) se $\operatorname{MCD}(a, b) = 1$.
\end{definizione}

\begin{esempio}
$3, 5$ sono coprimi; $11$ e $36$ sono coprimi.
\end{esempio}

\section{Minimo comune multiplo}

\begin{definizione}
Siano $a, b \in \mathbb{Z}$. Un elemento $m \in \mathbb{Z}$ è un \emph{minimo comune multiplo} di $a$ e $b$ se
\begin{enumerate}
  \item[i)] $a \mid m$ e $b \mid m$;
  \item[ii)] $\forall m' \in \mathbb{Z}$ tale che $a \mid m'$ e $b \mid m'$ si ha $m \mid m'$.
\end{enumerate}
\end{definizione}

Se $m$ è un minimo comune multiplo di $a$ e $b$, anche $-m$ lo è; indichiamo con $\operatorname{mcm}(a, b)$ quello \emph{positivo}.

\begin{esempio}
$\operatorname{mcm}(4, 6) = 12$, \quad $\operatorname{mcm}(-7, 9) = 63$.
\end{esempio}

\begin{proprieta}
$\operatorname{mcm}(a, b) = \operatorname{mcm}(b, a) = \operatorname{mcm}(|a|, |b|)$ e $\operatorname{mcm}(0, b) = 0$.
\end{proprieta}

\begin{proprieta}[Formula]
Per $a, b$ non entrambi nulli:
\[
  \operatorname{MCD}(a, b)\cdot\operatorname{mcm}(a, b) = |a|\cdot|b| .
\]
\end{proprieta}

\section{Algoritmo euclideo delle divisioni successive}

\begin{teorema}[Algoritmo di Euclide]
Siano $a, b \in \mathbb{Z}$, $a \ge b > 0$. Consideriamo queste uguaglianze (divisioni con resto):
\begin{align*}
  a &= q_1 b + r_1, & 0 &< r_1 < b,\\
  b &= q_2 r_1 + r_2, & 0 &< r_2 < r_1,\\
  r_1 &= q_3 r_2 + r_3, & 0 &< r_3 < r_2,\\
  &\ \,\vdots\\
  r_i &= q_{i+2}\, r_{i+1} + r_{i+2}, & 0 &< r_{i+2} < r_{i+1},\\
  &\ \,\vdots\\
  r_{n-3} &= q_{n-1}\, r_{n-2} + r_{n-1}, & 0 &< r_{n-1} < r_{n-2},\\
  r_{n-2} &= q_n\, r_{n-1} + r_n, & 0 &< r_n < r_{n-1},\\
  r_{n-1} &= q_{n+1}\, r_n .
\end{align*}
Allora $r_n = \operatorname{MCD}(a, b)$ (l'ultimo resto non nullo).
\end{teorema}

\begin{proof}
L'algoritmo termina in un numero finito di passi perché $0 < r_n < r_{n-1} < \dots < r_1 < b$.

\emph{Condizione i)} del MCD: dall'ultima uguaglianza $r_{n-1} = q_{n+1} r_n$, cioè $r_n \mid r_{n-1}$. Poi
\[
  r_{n-2} = r_{n-1} q_n + r_n = q_{n+1} r_n q_n + r_n \ \Rightarrow\ r_n \mid r_n(q_{n+1} q_n + 1) = r_{n-2}.
\]
Per il lemma, da $r_{n-3} = r_{n-2} q_{n-1} + r_{n-1}$ segue $r_n \mid r_{n-3}$. Risalendo tutte le uguaglianze si ottiene $r_n \mid r_i$ per ogni resto. Dalla seconda uguaglianza $b = r_1 q_2 + r_2$ e dal lemma segue $r_n \mid b$; dalla prima $a = q_1 b + r_1$ segue $r_n \mid a$. Questa è la condizione i) del MCD.

\emph{Condizione ii)}: sia $d$ un altro divisore comune di $a$ e $b$. Da $d \mid a = q_1 b + r_1$ e $d \mid b$ segue, per il lemma, $d \mid r_1$; poi $d \mid b = r_1 q_2 + r_2$ e $d \mid r_1$ implicano $d \mid r_2$; $d \mid r_1 = r_2 q_3 + r_3$ implica $d \mid r_3$, e così via. Quindi $d \mid r_i$ per ogni $i$, in particolare $d \mid r_{n-2}$ e $d \mid r_{n-1}$, da cui, per il lemma, $d \mid r_n = r_{n-2} - q_n r_{n-1}$. Cioè $d \mid r_n$: è la condizione ii) del MCD.
\end{proof}

\begin{esempio}[$\operatorname{MCD}(824, 376)$]
\begin{align*}
  824 &= 376\cdot 2 + 72,\\
  376 &= 72\cdot 5 + 16,\\
  72 &= 16\cdot 4 + 8,\\
  16 &= 8\cdot 2 + 0 .
\end{align*}
Dunque $\operatorname{MCD}(824, 376) = 8$ (l'ultimo resto non nullo).
\end{esempio}

\end{document}
